\SourceLeaf{452}%
\begin{center}
\begin{tabular}{@{}l p{.54\linewidth} r@{}}
\(a.\)&Die Länge von 1505 Messstangen ohne Correction \(.\ .\ .\ .\ .\ .\ .\ .\ .\)&\(3010{,}00000\) Toisen\\
\(b.\)&Summe der mit den Glaskeilen gemessenen Intervalle und der in Betracht kommenden ganzen und halben Durchmesser der Ablöthungs-Cylinder \(.\)&\(+3{,}58389\) T.\\
\(c.\)&Länge der Ergänzungsstange \(.\ .\ .\ .\ .\ .\ .\ .\ .\ .\ .\ .\ .\ .\ .\)&\(+1{,}22106\) T.\\
\(d.\)&Correction wegen Neigung der Ablöthungs-Cylinder gegen die Lothlinie \(.\)&\(-0{,}00008\) T.\\
\(e.\)&Correction wegen Abweichung der Stangen vom Alignement \(.\ .\ .\ .\ .\)&\(-0{,}00051\) T.\\
\(f.\)&Correction wegen fehlerhafter Längen der Messstangen \(.\ .\ .\ .\ .\ .\ .\)&\(-0{,}10245\) T.\\
\(g.\)&Correction wegen Abweichung der Temperatur der Messstangen von \(13^{0}\) R.&\(-0{,}19906\) T.\\
\(h.\)&Reduction auf die Oberfläche des Meeres \(.\ .\ .\ .\ .\ .\ .\ .\ .\ .\ .\)&\(-0{,}02264\) T.\\\cline{3-3}
&Länge der Braacker-Basis nach der neuen Berechnung \(.\ .\ .\ .\ .\)&\(=3014{,}48021\) Toisen
\end{tabular}
\end{center}

Es findet sich aber auch in dieser Berechnung ein schwacher Punkt, nämlich die sub \(g\) angeführte Correction wegen der Temperatur der Messstangen. Eine mit Abbildungen versehene Beschreibung des bei der Basismessung angewandten Apparats hat \textsc{Schumacher} in der Schrift: ‘Schreiben an Dr. \textsc{Olbers} in Bremen etc. etc.’, Altona 1821’ veröffentlicht, und man wird daraus ersehen, dass die Temperaturen nicht durch Metallthermometer, sondern durch gewöhnliche, eingelegte Thermometer bestimmt sind. Dies ist nun an und für sich ein misslicher Umstand, aber viel schlimmer stellt sich die Sache, da die Ausdehnbarkeit der Stangen nur aus einigen im Felde vorgenommenen Messungen der Stangenlängen am Abend und am Morgen abgeleitet wird. Es kann aber diesem Uebel abgeholfen werden. Im Jahre 1853 wurde nemlich die Stange No. IV. des \textsc{Schumacher}schen Basisapparats nach Pulkowa gebracht, um direct mit den dort gesammelten Etalons verglichen zu werden. Bei dieser Gelegenheit wurde nun auch die Ausdehnung dieser Stange für \(100^{0}\) erhalten, und wenn man den von \textsc{Struve} (Siehe ‘Arc du méridien entre le Danube et la mer glaciale’ pag. 51) angegebenen Werth der Ausdehnungscoefficienten berechnet, dann erhält man für die Correction sub \(g\): \(-0{,}22812\) statt \(-0{,}19906\).

Mit dieser Berichtigung, welche auch von Professor \textsc{Peters} adoptirt wird, findet man dann die Länge der \textit{Braacker Basis}:
\[
=3014{,}451\ \text{Toisen},
\]
und dieser Werth muss als der \textit{definitive} betrachtet werden. Ich füge nun hinzu, dass die Angabe dieser Toisen auf der Vergleichung mit der Pulkowaer \textsc{Fortin} beruhe; da diese aber mit der \textsc{Bessel}’schen Toise bis auf eine verschwindende Kleinigkeit übereinstimmt, kann die Länge auch füglich als in \textsc{Bessel}schen Toisen ausgedrückt angesehen werden.’

Obiges Coordinaten-Verzeichniss ergiebt für die Länge der Basis \(5875{,}3614\) der dort angewandten Einheiten oder \(3014{,}5757\) Toisen bei einem Erdmeridian von \(360\times57009{,}746\) Toisen.

\textsc{Eduard Schmidt. Gauss} an \textsc{Schumacher}: Göttingen 1830 April 30. ‘Um Ihr Vertrauen zu \textsc{Schmidt}’s Rechnung zu vergrössern, bemerke ich, dass er die zwei Hauptelemente der Erddimensionen viermal berechnet hat. — — Das Resultat (IV) ist mir von ihm handschriftlich mitgetheilt und dasselbe was meinen neuen Hülfstafeln zum Grunde liegt, nemlich Abplattung \(\frac{1}{297{,}732}\); \(\frac{\text{Erd-Quadrant}}{90}=57008^{\mathrm{T}}551\).’

\textsc{Bessel}. ‘Ueber einen Fehler in der Berechnung der französischen Gradmessung und seinen Einfluss auf die Bestimmung der Figur der Erde.’ Astronomische Nachrichten Nr. 438 Band 19. Seite 116. 1841 December 2. ‘Mittlere Grad des Meridians \(=57013{,}109\) Toisen, halbe grosse Axe \(a=3272077{,}14\) Toisen, halbe kleine Axe \(b=3261139{,}33\) Toisen, \(a:b=299{,}1528:298{,}1528\)’

Bei der Anwendung der in obigen Verzeichnissen angegebenen Coordinaten hat man diese also vorläufig, ehe die Basis und die Verbindungsdreiecke bis \textit{Hamburg — Hohenhorn} von Neuem gemessen sind, mit folgendem Correctionsfactor zu multipliciren:
